\SixthArticle{33}
Wie wir im vorigen Artikel $\Psi\left(-\tfrac12\right)$ auf $\Psi0$ und einen Logarithmus zurückgeführt haben, so werden wir allgemein $\Psi\left(-\frac mn\right)$, wobei $m,n$ ganze Zahlen bezeichnen, von denen $m$ die kleinere ist, auf $\Psi0$ und Logarithmen zurückführen. Setzen wir $\frac{2\pi}{n}=\omega$, und sei $\varphi$ gleich einem der Winkel $\omega,2\omega,3\omega\ .\ .\ .\ .\ (n-1)\omega$; dann ist $1=\cos n\varphi=\cos2n\varphi=\cos3n\varphi$ usw., $\cos\varphi=\cos(n+1)\varphi=\cos(n+2)\varphi$ usw., $\cos2\varphi=\cos(n+2)\varphi$ usw. und außerdem $\cos\varphi+\cos2\varphi+\cos3\varphi+\text{usw.}+\cos(n-1)\varphi+1=0$. Wir haben also
\[
\begin{aligned}
&\cos\varphi\cdot\Psi\frac{1-n}{n}=-n\cos\varphi+\cos\varphi\cdot\log2\\[4pt]
&-\frac{n}{n+1}\cos(n+1)\varphi+\cos\varphi\cdot\log\tfrac32-\text{usw.}
\end{aligned}
\]
\[
\begin{aligned}
&\cos2\varphi\cdot\Psi\frac{2-n}{n}=-\frac n2\cos2\varphi+\cos2\varphi\cdot\log2\\[4pt]
&-\frac{n}{n+2}\cos(n+2)\varphi+\cos2\varphi\cdot\log\tfrac32-\text{usw.}
\end{aligned}
\]
\[
\begin{aligned}
&\cos3\varphi\cdot\Psi\frac{3-n}{n}=-\frac n3\cos3\varphi+\cos3\varphi\cdot\log2\\[4pt]
&-\frac{n}{n+3}\cos(n+3)\varphi+\cos3\varphi\cdot\log\tfrac32-\text{usw.}
\end{aligned}
\]
usw. bis zu
\SourceLeaf{165}
\[
\begin{aligned}
&\cos(n-1)\varphi\cdot\Psi\left(-\frac1n\right)
\\[4pt]
&=-\frac{n}{n-1}\cos(n-1)\varphi+\cos(n-1)\varphi\cdot\log2\\[4pt]
&-\frac{n}{2n-1}\cos(2n-1)\varphi
+\cos(n-1)\varphi\cdot\log\tfrac32-\text{usw.}
\end{aligned}
\]
\[
\Psi0=-\frac nn\cos n\varphi+\log2-\frac{n}{2n}\cos2n\varphi+\log\tfrac32-\text{usw.}
\]
und durch \textit{Summation}
\[
\begin{aligned}
&\cos\varphi\cdot\Psi\frac{1-n}{n}+\cos2\varphi\cdot\Psi\frac{2-n}{n}\\[4pt]
&+\cos3\varphi\cdot\Psi\frac{3-n}{n}+\text{usw.}\\[4pt]
&+\cos(n-1)\varphi\cdot\Psi\left(-\frac1n\right)+\Psi0
\\[4pt]
&=-n(\cos\varphi+\tfrac12\cos2\varphi+\tfrac13\cos3\varphi+\tfrac14\cos4\varphi+\text{usw. ins Unendliche})
\end{aligned}
\]
Für einen Wert von $x$, der nicht größer als die Einheit ist, gilt aber allgemein
\[
\begin{aligned}
&\log(1-2x\cos\varphi+xx)\\[4pt]
&=-2(x\cos\varphi+\tfrac12xx\cos2\varphi+\tfrac13x^3\cos3\varphi+\text{usw.})
\end{aligned}
\]
Diese Reihe ergibt sich leicht aus der Entwicklung von $\log(1-rx)+\log\left(1-\frac xr\right)$, wobei $r$ die Größe $\cos\varphi+\BareRoot{-1}\cdot\sin\varphi$ bezeichnet. Hierdurch geht die vorige Gleichung über in
\[
\begin{aligned}
&\cos\varphi\cdot\Psi\frac{1-n}{n}+\cos2\varphi\cdot\Psi\frac{2-n}{n}\\[4pt]
&+\cos3\varphi\cdot\Psi\frac{3-n}{n}+\text{usw.}\\[4pt]
&+\cos(n-1)\varphi\cdot\Psi\left(-\frac1n\right)
\\[4pt]
&=-\Psi0+\tfrac12n\log(2-2\cos\varphi)
\end{aligned}
\tag*{[72]}\]
In dieser Gleichung setze man der Reihe nach $\varphi=\omega$, $\varphi=2\omega$, $\varphi=3\omega$ usw. bis $\varphi=(n-1)\omega$, multipliziere diese Gleichungen in derselben Reihenfolge mit $\cos m\omega$, $\cos2m\omega$, $\cos3m\omega$ usw. bis $\cos(n-1)m\omega$ und füge zur Summe der Produkte Gleichung 71 hinzu:
\[
\begin{aligned}
&\Psi\frac{1-n}{n}+\Psi\frac{2-n}{n}+\Psi\frac{3-n}{n}+\text{usw.}+\Psi\left(-\frac1n\right)\\[4pt]
&=(n-1)\Psi0-n\log n
\end{aligned}
\]
Berücksichtigt man nun, dass
\[
\begin{aligned}
&1+\cos m\omega\cdot\cos k\omega+\cos2m\omega\cdot\cos2k\omega\\[4pt]
&+\cos3m\omega\cdot\cos3k\omega
+\text{usw.}+\cos(n-1)m\omega\cdot\cos(n-1)k\omega=0
\end{aligned}
\]
gilt, wobei $k$ eine der Zahlen $1,2,3\ .\ .\ .\ .\ (n-1)$ mit Ausnahme der beiden $m$ und $n-m$ bezeichnet, für die jene Summe $=\tfrac12n$ wird, so ist ersichtlich, dass aus der Summation jener Gleichungen nach Division durch $\frac n2$ folgt:
\[
\begin{aligned}
&\Psi\left(-\frac mn\right)+\Psi\left(-\frac{n-m}{m}\right)=
\\[4pt]
&2\Psi0-2\log n+\cos m\omega\cdot\log(2-2\cos\omega)\\[4pt]
&+\cos2m\omega\cdot\log(2-2\cos2\omega)
\\[4pt]
&+\cos3m\omega\cdot\log(2-2\cos3\omega)+\text{usw.}\\[4pt]
&+\cos(n-1)m\omega\cdot\log(2-2\cos(n-1)\omega)
\end{aligned}
\tag*{[73]}\]
Offenbar ist das letzte Glied dieser Gleichung $=\cos m\omega\cdot\log(2-2\cos\omega)$, das\SourceLeaf{166} vorletzte $=\cos 2m\omega.\log(2-2\cos 2\omega)$ usw., sodass stets je zwei Glieder gleich sind, mit Ausnahme, wenn $n$ gerade ist, des einzelnen Gliedes $\cos\frac{n}{2}.m\omega\log(2-2\cos\frac{n}{2}\omega)$; dieses wird $=+2\log 2$ für gerades $m$ oder $=-2\log 2$ für \mbox{ungerades $m$}. Verbinden wir nun mit Gleichung 73 die Gleichung
\[
\Psi\left(-\frac{m}{n}\right)-\Psi\left(-\frac{n-m}{n}\right)=\pi\operatorname{cotang}\frac{m}{n}\pi
\]
so haben wir für einen ungeraden Wert von $n$, sofern $m$ eine positive ganze Zahl kleiner als $n$ ist,
\[
\begin{aligned}
&{}\text{[74]}\qquad \Psi\left(-\frac{m}{n}\right)
\\[4pt]
&=\Psi 0+\tfrac12\pi\operatorname{cotang}\frac{m\pi}{n}-\log n\\[4pt]
&+\cos\frac{2m\pi}{n}.\log\left(2-2\cos\frac{2\pi}{n}\right)
\\[4pt]
&+\cos\frac{4m\pi}{n}.\log\left(2-2\cos\frac{4\pi}{n}\right)\\[4pt]
&+\cos\frac{6m\pi}{n}.\log\left(2-2\cos\frac{6\pi}{n}\right)+\text{usw.}
\\[4pt]
&+\cos\frac{(n-1)m\pi}{n}.\log\left(2-2\cos\frac{(n-1)\pi}{n}\right)
\end{aligned}
\]
Für einen geraden Wert von $n$ dagegen
\[
\begin{aligned}
&{}\text{[75]}\qquad \Psi\left(-\frac{m}{n}\right)
\\[4pt]
&=\Psi 0+\tfrac12\pi\operatorname{cotang}\frac{m\pi}{n}-\log n\\[4pt]
&+\cos\frac{2m\pi}{n}\log\left(2-2\cos\frac{2\pi}{n}\right)
\\[4pt]
&+\cos\frac{4m\pi}{n}\log\left(2-2\cos\frac{4\pi}{n}\right)+\text{usw.}\\[4pt]
&+\cos\frac{(n-2)m\pi}{n}\log\left(2-2\cos\frac{(n-2)\pi}{n}\right)
\pm\log 2
\end{aligned}
\]
wobei das obere Zeichen für gerades $m$, das untere für ungerades gilt. So ergibt sich beispielsweise
\[
\begin{aligned}
&\Psi\left(-\tfrac14\right)=\Psi 0+\tfrac12\pi-3\log 2, 
\\[4pt]
&\Psi\left(-\tfrac34\right)=\Psi 0-\tfrac12\pi-3\log 2
\\[4pt]
&\Psi\left(-\tfrac13\right)=\Psi 0+\tfrac12\pi\BareRoot{\tfrac13}-\tfrac32\log 3, 
\\[4pt]
&\Psi\left(-\tfrac23\right)=\Psi 0-\tfrac12\pi\BareRoot{\tfrac13}-\tfrac22\log 3
\end{aligned}
\]
Verbindet man im Übrigen diese Gleichungen mit Gleichung 64, so zeigt sich von selbst, dass sich $\Psi z$ allgemein für \emph{jeden rationalen Wert} von $z$, sei er positiv oder negativ, durch $\Psi 0$ und Logarithmen bestimmen lässt; dieser Satz ist gewiss höchst bemerkenswert.

\SixthArticle{34}
Da nach Artikel 28 $\Pi\lambda$ der Wert des Integrals $\int y^\lambda e^{-y}\mathrm{d}y$ von $y=0$ bis $y=\infty$ ist, sofern $\lambda+1$ eine positive Größe ist, ergibt sich durch Differentiation nach $\lambda$
\[
\frac{\mathrm{d}\Pi\lambda}{\mathrm{d}\lambda}
=\frac{\mathrm{d}\int y^\lambda e^{-y}\mathrm{d}y}{\mathrm{d}\lambda}
=\int y^\lambda e^{-y}\log y\,\mathrm{d}y
\]
oder
\[
\begin{aligned}
&[76]\qquad \Pi\lambda.\Psi\lambda=\int y^\lambda e^{-y}\log y.\mathrm{d}y,\\[4pt]
&\quad\text{von }y=0\text{ bis }y=\infty
\end{aligned}
\]
\par\SourceLeaf{167}
{\raggedright\parindent=1.2em Setzt man allgemeiner $y=z^\alpha$, $\alpha\lambda+\alpha-1=\beta$, so ist der Wert des Integrals $\int z^\beta e^{-z^\alpha}\log z\,\mathrm{d}z$ von $z=0$ bis $z=\infty$ gleich\par}
\[
\begin{aligned}
&=\frac{1}{\alpha\alpha}\Pi\left(\frac{\beta+1}{\alpha}-1\right).\Psi\left(\frac{\beta+1}{\alpha}-1\right)
\\[4pt]
&=\frac{1}{\alpha(\beta+1)}\Pi\frac{\beta+1}{\alpha}.\Psi\frac{\beta+1}{\alpha}
-\frac{1}{(\beta+1)^2}\Pi\frac{\beta+1}{\alpha}
\end{aligned}
\]
sofern $\beta+1$ und $\alpha$ zugleich positive Größen sind; oder er ist derselben Größe mit entgegengesetztem Zeichen gleich, wenn sowohl $\beta+1$ als auch $\alpha$ negativ sind.

